
\section{\BookRecordIndex{T0127}分支接触}\label{sec:E-6}

本节完成定理~\ref{thm:E-1-1} 证明的第二步：精确分析两条不同的 \BookRecordIndex{T0003}$J$-全纯曲线在交点附近的距离。唯一非平凡的情形，是两曲线在交点处有相同的切空间与奇点阶数。由于问题是局部的，将 $M$ 识别为 $\mathbb C^n$，并假定曲线具有命题~\ref{prop:E-5-10} 给出的形式。

\begin{proposition}\label{prop:E-6-1}
设 $J$ 是 $\mathbb C^n$ 上在原点为标准结构的\BookRecordIndex{T0088}殆复结构，$k$ 为正整数。对 $j=1,2$，设 $U_j\subset\mathbb C$ 为原点的开邻域，$u_j:\overline U_j\to\mathbb C^n$ 是 \BookRecordIndex{T0003}$J$-全纯曲线，$\varphi_j:\mathbb D\to\overline U_j$ 是 $C^2$ 微分同胚，满足 $\varphi_j-\mathrm{id}\in\mathcal O_2$；当 $k=1$ 时还要求 $\varphi_j$ 光滑，并且
\[
  u_j\circ\varphi_j(z)=(z^k,F_j(z)),\qquad j=1,2,
\]
其中 $F_j:\mathbb D\to\mathbb C^{n-1}$ 属于 $\mathcal O_{k+1}$。若 $F_2-F_1$ 在零点的任何邻域内均不恒为零，则存在整数 $m>k$ 及非零向量 $a\in\mathbb C^{n-1}$，使
\[
  F_2(z)-F_1(z)=z^ma+R(z),\qquad R\in\mathcal O_{m+1}.
\]
\end{proposition}

证明思路如下。若重参数化 $\varphi_j$ 是全纯的，便可对差 $F_1-F_2$ 应用 \BookRecordIndex{T0024}Hartman--Wintner 定理，得到某个整数 $m>k$ 及非零的 $m$ 次齐次多项式 $h$，使 $F_1-F_2-h\in\mathcal O_{m+1}$。引理~\ref{lem:E-5-9} 随后可\BookRecordIndex{T0326}推出 $h$ 全纯，故为 $z^ma$ 的形式。但这里仅知复合映射 $u_j\circ\varphi_j$ 在每一点都有复切空间，须做更多工作才能验证 $F_1-F_2$ 满足 \BookRecordIndex{T0024}Hartman--Wintner 定理的假设。首先要在 $\mathbb C^n$ 的原点附近选取合适的\emph{\BookRecordIndex{T0400}正规坐标}，再考察此坐标下 $F_j$ 的性质。

\paragraph{\BookRecordIndex{T0400}正规坐标。}
下述命题构造与给定一维复子空间 $\Pi_0\subset T_{x_0}M$ 相切的嵌入 $J$-全纯圆盘，并将它扩充为一个以复余维一空间为参数的嵌入 $J$-全纯圆盘族。应用时，$x_0$ 是 \BookRecordIndex{T0003}$J$-全纯曲线的奇点，$\Pi_0$ 是该曲线的切空间；结果是在奇点附近得到一个有用的坐标图。

\begin{proposition}\label{prop:E-6-2}
设 $(M,J)$ 是光滑的殆复 $2n$-流形，$x_0\in M$，$\Pi_0\subset T_{x_0}M$ 是复直线。则存在光滑坐标图 $\Psi:(M,x_0)\to(\mathbb C^n,0)$，使
\[
  \Psi(x_0)=0,\qquad d\Psi(x_0)\Pi_0=\mathbb C\times\{0\},
  \qquad
  \Psi_*J=\begin{pmatrix}i&J_3\\0&J_4\end{pmatrix}
  \in\mathrm{GL}_{\mathbb R}(\mathbb C\times\mathbb C^{n-1}),
\]
其中 $J_3(0)=0$，$J_4(0)$ 是 $\mathbb C^{n-1}$ 上的标准复结构 $i$。
\end{proposition}

\begin{proof}
不妨设 $M=\mathbb C^n$、$x_0=0$，且 $J$ 是满足 $J(0)=i$ 的\BookRecordIndex{T0088}殆复结构。固定 $p>2$，定义
\[
  \mathcal X:=\left\{u\in W^{2,p}(\mathbb D,\mathbb C^n)\mid
    u(e^{i\theta})\in e^{3i\theta/2}\mathbb R^n\right\},
  \qquad
  \mathcal Y:=W^{1,p}(\mathbb D,\mathbb C^n)
     \times\mathbb C^n\times\mathbb C^n.
\]
定义 $\mathcal F:\mathcal X\to\mathcal Y$：
\[
  \mathcal F(u):=\left(\tfrac12(\partial_su+J(u)\partial_tu),
                        u(0),\partial_su(0)\right).
\]
于是对 $\xi\in\mathcal X$，有 $d\mathcal F(0)\xi=(\bar\partial\xi,\xi(0),\partial_s\xi(0))$。与定理~\ref{thm:E-3-1} 证明第 3 步相同的论证表明，$d\mathcal F(0)$ 是 Banach 空间同构。故由逆函数定理，$\mathcal F$ 在 $\mathcal Y$ 中零点的某开邻域 $\mathcal U$ 上有光滑逆映射，且将零映为零。

取 $\Pi_0$ 在 $\mathbb C^n$ 中的复补空间 $\Pi$，取 $\xi_0\in\Pi_0\setminus\{0\}$，并取 $\Pi$ 中原点的开邻域 $W$，使任意 $x\in W$ 均有 $(0,x,\xi_0)\in\mathcal U$。\sourcefootnote{原书实体页 694 写 $(0,\xi_0,x)\in\mathcal U$，但 $\mathcal F(u)=(\bar\partial_Ju,u(0),\partial_su(0))$，下一式又取 $\mathcal F^{-1}(0,x,\xi_0)$；故按此唯一参数顺序补正。}定义
\[
  \Phi:\mathbb D\times W\to\mathbb C^n,\qquad
  \Phi(z,x):=\mathcal F^{-1}(0,x,\xi_0)(z).
\]
对 $x\in W$，令 $u_x(z):=\Phi(z,x)$。则 $u_0(z)=z\xi_0+z^2\bar\xi_0$；$u_x$ 是 $u_0$ 附近满足 $u_x(0)=x$、$\partial_su_x(0)=\xi_0$ 及 $u_x(e^{i\theta})\in e^{3i\theta/2}\mathbb R^n$ 的唯一 $J$-全纯圆盘。\sourcefootnote{原书实体页 695 由 $u_0=\mathcal F^{-1}(0,0,\xi_0)$ 直接写出 $u_0(z)=z\xi_0+z^2\bar\xi_0$。该多项式满足标准复结构的全纯方程，但在此仅假定 $J(0)=i$，一般非恒定 $J$ 下未必是 $J$-全纯解；这里保留原断言并标出疑点。}由\BookRecordIndex{T0329}椭圆正则性，$u_x$ 光滑依赖于 $x$，故 $\Phi$ 光滑。又有
\[
  d\Phi(0)(\zeta,\xi)=\zeta\xi_0+\xi.
\]
于是 $\Phi(0)=0$ 且 $d\Phi(0)\in\mathrm{GL}(n,\mathbb C)$。由逆函数定理，$\Phi$ 在 $\mathbb D\times\Pi$ 中零点的足够小邻域上限制为到 $\mathbb C^n$ 中零点邻域的微分同胚。将 $\Pi$ 识别为 $\mathbb C^{n-1}$ 后，$\Psi:=\Phi^{-1}$ 即满足命题要求。
\end{proof}

\paragraph{具有复切空间的映射。}
现在回到命题~\ref{prop:E-6-1}。首先要证明，尽管重参数化
\[
  u\circ\varphi(z)=(z^k,F(z))
\]
未必是 $J$-全纯的，$F$ 在\BookRecordIndex{T0400}正规坐标下仍满足椭圆方程。证明只需直接运算，但代数整理较多。记
\[
  \{\Phi,F\}:=\partial_s\Phi\cdot\partial_tF
                    -\partial_t\Phi\cdot\partial_sF
\]
为函数 $\Phi:\mathbb D\to\operatorname{Hom}_{\mathbb R}(\mathbb C^k,\mathbb C^\ell)$ 与 $F:\mathbb D\to\mathbb C^k$ 的 Poisson 括号。

\begin{lemma}\label{lem:E-6-3}
设
\begin{equation}\label{eq:E-6-1}
  J:=\begin{pmatrix}i&J_3\\0&J_4\end{pmatrix}
  \in\operatorname{End}_{\mathbb R}(\mathbb C\oplus\mathbb C^{n-1})
\end{equation}
是 $\mathbb C^n$ 上的 $C^1$ \BookRecordIndex{T0088}殆复结构，$u:\mathbb D\to\mathbb C^n=\mathbb C\times\mathbb C^{n-1}$ 是形如 $u(z)=(f(z),F(z))$ 的 $C^2$ 映射。设 $f:\mathbb D\to\mathbb C$ 为非常值全纯函数，记 $Y:=\{z\in\mathbb D\mid f'(z)=0\}$，并假定 $u$ 在每一点 $z\in\mathbb D\setminus Y$ 都有一维 $J$-复切空间。定义 $A,B,C:\mathbb D\setminus Y\to\mathbb C$：
\begin{equation}\label{eq:E-6-2}
  A:=\frac1{f'}J_3(u)\partial_sF,\qquad
  B:=\frac1{f'}J_3(u)\partial_tF,\qquad
  C:=\left\{\frac1{f'}J_3(u),F\right\}.
\end{equation}
则 $F$ 在 $\mathbb D\setminus Y$ 上满足
\begin{align}
  \partial_sF+J_4(u)\cdot\partial_tF
    &=\operatorname{Re}(B)\partial_sF
       +\operatorname{Im}(B)\partial_tF,\label{eq:E-6-3}\\
  J_4(u)\cdot\partial_sF-\partial_tF
    &=\operatorname{Re}(A)\partial_sF
       +\operatorname{Im}(A)\partial_tF,\label{eq:E-6-4}
\end{align}
以及
\begin{align}
  \Delta F-\operatorname{Re}(B)\partial_s^2F
  +\operatorname{Im}(A)\partial_t^2F
  +\operatorname{Re}(A+iB)\partial_s\partial_tF
  &=\operatorname{Re}(C)\partial_sF
     +\operatorname{Im}(C)\partial_tF
     -\{J_4(u),F\}.\label{eq:E-6-5}
\end{align}
\end{lemma}

\begin{proof}
由于 $u$ 具有 $J$-复切空间，存在实值函数 $a,b,c,d:\mathbb D\setminus Y\to\mathbb R$，满足
\begin{align}
  J(u)\cdot(\partial_sf,\partial_sF)
    &=a\cdot(\partial_sf,\partial_sF)
      +b\cdot(\partial_tf,\partial_tF),\label{eq:E-6-6}\\
  J(u)\cdot(\partial_tf,\partial_tF)
    &=c\cdot(\partial_sf,\partial_sF)
      +d\cdot(\partial_tf,\partial_tF).\label{eq:E-6-7}
\end{align}
因 $\partial_sf=f'$、$\partial_tf=if'$，取上两式的首分量，得到
\[
  if'+J_3(f,F)\partial_sF=af'+ibf',\qquad
  -f'+J_3(f,F)\partial_tF=cf'+idf'.
\]
除以 $f'$ 并用式~\eqref{eq:E-6-2}，可得
\[
  a=\operatorname{Re}(A),\quad b=1+\operatorname{Im}(A),
  \qquad c=-1+\operatorname{Re}(B),\quad d=\operatorname{Im}(B).
\]
取式~\eqref{eq:E-6-6}、\eqref{eq:E-6-7} 的第二分量，即得式~\eqref{eq:E-6-3}、\eqref{eq:E-6-4}。计算 $\partial_s(\partial_sF+J_4\partial_tF)+\partial_t(\partial_tF-J_4\partial_sF)$，再用 $C=\partial_sB-\partial_tA$，即得式~\eqref{eq:E-6-5}。
\end{proof}

下面考察引理~\ref{lem:E-6-3} 中两映射 $u_j=(f_j,F_j)$ 的情形，其中 $f_1(z)=f_2(z)=z^k$，并研究差 $H:=F_2-F_1$\sourcefootnote{原书实体页 696 的引导段及实体页 698 的命题证明开头曾把 $H$ 写作 $F_1-F_2$，但紧接的引理 E.6.4、命题 E.6.1 的陈述及后续差分计算均取 $H=F_2-F_1$；据这些固定定义统一记号。}。关键是下述式~\eqref{eq:E-6-9}：它表明，对适当的实函数 $a,b,c:\mathbb D\to\mathbb R$，$H$ 满足 \BookRecordIndex{T0024}Hartman--Wintner 定理的假设。

\begin{lemma}\label{lem:E-6-4}
设 $J$ 如引理~\ref{lem:E-6-3}，且 $J(0)=i$。设 $k$ 为正整数，$F_1,F_2:\mathbb D\to\mathbb C^{n-1}$ 均属于 $\mathcal O_{k+1}$。记
\[
  u_j(z):=(z^k,F_j(z)),\qquad H:=F_2-F_1,
\]
并假定 $u_j$ 在每一点 $z\in\mathbb D\setminus\{0\}$ 均有 $J$-复切空间。则：
\begin{enumerate}
\item[(i)] 映射 $J_3(u_2)-J_3(u_1)$ 与 $J_4(u_2)-J_4(u_1)$ 可写为
\[
  J_3(u_2)(z)-J_3(u_1)(z)=\alpha(z)\cdot H(z),\qquad
  J_4(u_2)(z)-J_4(u_1)(z)=\beta(z)\cdot H(z),
\]
其中 $\alpha:\mathbb D\to\operatorname{Hom}_{\mathbb R}(\mathbb C^{n-1},
\operatorname{Hom}_{\mathbb R}(\mathbb C^{n-1},\mathbb C))$、
$\beta:\mathbb D\to\operatorname{Hom}_{\mathbb R}(\mathbb C^{n-1},
\operatorname{End}_{\mathbb R}(\mathbb C^{n-1}))$ 均属于 $\mathcal O_0$。
\item[(ii)] 映射 $\bar\partial H:=\tfrac12(\partial_sH+J_0\partial_tH)$ 具有形式
\begin{equation}\label{eq:E-6-8}
  \bar\partial H(z)=\varepsilon_1(z)\cdot H(z)
      +\varepsilon_2(z)\cdot\partial_sH(z)
      +\varepsilon_3(z)\cdot\partial_tH(z),
\end{equation}
其中 $\varepsilon_1,\varepsilon_2,\varepsilon_3:\mathbb D\to
\operatorname{End}_{\mathbb R}(\mathbb C^{n-1})$ 属于 $\mathcal O_k$。
\item[(iii)] 存在常数 $C>0$ 与实函数 $a,b,c:\mathbb D\to\mathbb R$，其中 $a-1,b,c-1\in\mathcal O_{k+1}$，满足 $a(0)=c(0)=1$、$b(0)=0$，并在零点附近满足\sourcefootnote{原书实体页 696 把 $a,b,c$ 本身均写作属于 $\mathcal O_{k+1}$，却同时要求 $a(0)=c(0)=1$；二者不\BookRecordIndex{T0348}相容。由本引理证明中的 $a=1-\operatorname{Re}B_2$、$2b=\operatorname{Re}(A_2+iB_2)$、$c=1+\operatorname{Im}A_2$ 及 $A_2,B_2\in\mathcal O_{k+1}$，唯一补正是 $a-1,b,c-1\in\mathcal O_{k+1}$。}
\begin{equation}\label{eq:E-6-9}
  |a\partial_s^2H+2b\partial_s\partial_tH+c\partial_t^2H|
  \le C\bigl(|H|+|\partial_sH|+|\partial_tH|\bigr).
\end{equation}
\end{enumerate}
\end{lemma}

\begin{proof}
对 $j=1,2$，按式~\eqref{eq:E-6-2} 定义 $A_j,B_j,C_j:\mathbb D\setminus\{0\}\to\mathbb C$，其中 $f(z)=z^k$，且以 $(u_j,F_j)$ 代替 $(u,F)$。则 $A_j,B_j\in\mathcal O_{k+1}$，$C_j\in\mathcal O_k$。事实上，$1/f'\in\mathcal O_{1-k}$、$\partial_sF_j\in\mathcal O_k$、$J_3(0)=0$，由引理~\ref{lem:E-5-3} 的复合性质及注~\ref{rem:E-5-5}，$J_3(f,F_j)\in\mathcal O_k$。故由乘积性质，$A_j=(1/f')J_3(f,F_j)\partial_sF_j\in\mathcal O_{k+1}$；$B_j$ 同理，进而 $C_j=\partial_sB_j-\partial_tA_j\in\mathcal O_k$。

先证 (i)。定义 $\alpha:\mathbb D\to\operatorname{Hom}_{\mathbb R}(\mathbb C^{n-1},\operatorname{Hom}_{\mathbb R}(\mathbb C^{n-1},\mathbb C))$：
\[
  \alpha(z)\xi:=\left(\int_0^1
       dJ_3\bigl(u_1(z)+\tau(0,H(z))\bigr)\,d\tau\right)
       \cdot(0,\xi),
\]
其中 $z\in\mathbb D$、$\xi\in\mathbb C^{n-1}$。由 $(0,H)=u_2-u_1$，
\[
  J_3(u_2)(z)-J_3(u_1)(z)=\alpha(z)\cdot H(z).
\]
这可与定理~2.3.2 证明中的 $J$ 公式比较。须证 $\alpha\in\mathcal O_0$。记
\[
  \alpha_\tau(z)\xi:=dJ_3\bigl(u_1(z)+\tau(0,H(z))\bigr)
      \cdot(0,\xi),\qquad
  \widehat\alpha\xi:=dJ_3(0)\cdot(0,\xi).
\]
由注~\ref{rem:E-5-5} 的复合性质，对每个 $\tau$ 均有 $\alpha_\tau-\widehat\alpha\in\mathcal O_k$，而 $\mathcal O_k$ 定义中的常数可独立于 $\tau$ 选取。因此\sourcefootnote{原书实体页 697 此式积分上限印作 $\tau$，但 $\alpha$ 在紧邻定义为从 $0$ 到 $1$ 的积分；据定义改为上限 $1$。} $\alpha-\widehat\alpha=\int_0^1(\alpha_\tau-\widehat\alpha)\,d\tau\in\mathcal O_k$，从而 $\alpha\in\mathcal O_0$。关于 $\beta$ 的证明相同。

再证 (ii)。由 $\alpha$ 的定义，
\[
  B_2-B_1=\frac1{f'}\bigl((\alpha\cdot H)\cdot\partial_tF_1
                       +J_3(u_2)\cdot\partial_tH\bigr).
\]
因 $1/f'\in\mathcal O_{1-k}$、$\partial_tF_1\in\mathcal O_k$、$\alpha\in\mathcal O_0$，乘积性质给出
\[
  (B_2-B_1)(z)=\beta_1(z)\cdot H(z)
                     +\beta_2(z)\cdot\partial_tH(z),
\]
其中 $\beta_1,\beta_2:\mathbb D\to\operatorname{Hom}_{\mathbb R}(\mathbb C^{n-1},\mathbb C)$ 均属于 $\mathcal O_1$。由引理~\ref{lem:E-6-3}，$F_j$ 满足式~\eqref{eq:E-6-3}：
\[
  \partial_sF_j+J_4(u_j)\cdot\partial_tF_j
   =\operatorname{Re}(B_j)\partial_sF_j
     +\operatorname{Im}(B_j)\partial_tF_j.
\]
两式相减得
\begin{align*}
  2\bar\partial H+(J_4(u_2)-J_4(u_1))\partial_tF_1
   +(J_4(u_2)-i)\partial_tH
   &=\operatorname{Re}(B_2-B_1)\partial_sF_1
     +\operatorname{Re}(B_2)\partial_sH\\
   &\quad+\operatorname{Im}(B_2-B_1)\partial_tF_1
     +\operatorname{Im}(B_2)\partial_tH.
\end{align*}
于是
\begin{align*}
  2\bar\partial H+(\beta\cdot H)\partial_tF_1
     +(J_4(u_2)-i)\partial_tH
  &=\operatorname{Re}(\beta_1\cdot H+\beta_2\cdot\partial_tH)
       \partial_sF_1
       +\operatorname{Re}(B_2)\partial_sH\\
  &\quad+\operatorname{Im}(\beta_1\cdot H+\beta_2\cdot\partial_tH)
       \partial_tF_1+\operatorname{Im}(B_2)\partial_tH.
\end{align*}
由于 $J_4(u_2)-i\in\mathcal O_k$，这便证明了 (ii)。

最后证 (iii)。类似计算表明
\begin{align*}
  A_1(z)-A_2(z)&=\alpha_1(z)\cdot H(z)
                     +\alpha_2(z)\cdot\partial_sH(z),\\
  C_1(z)-C_2(z)&=\gamma_1(z)\cdot H(z)
                     +\gamma_2(z)\cdot\partial_sH(z)
                     +\gamma_3(z)\cdot\partial_tH(z),
\end{align*}
其中 $\alpha_1,\alpha_2\in\mathcal O_1$、$\gamma_1,\gamma_2,\gamma_3\in\mathcal O_0$，它们的值均属于 $\operatorname{Hom}_{\mathbb R}(\mathbb C^{n-1},\mathbb C)$。由引理~\ref{lem:E-6-3}，$F_1,F_2$ 满足式~\eqref{eq:E-6-5}，即对 $j=1,2$ 有
\begin{multline*}
  \Delta F_j-\operatorname{Re}(B_j)\partial_s^2F_j
   +\operatorname{Im}(A_j)\partial_t^2F_j
   +\operatorname{Re}(A_j+iB_j)\partial_s\partial_tF_j\\
  =\operatorname{Re}(C_j)\partial_sF_j
    +\operatorname{Im}(C_j)\partial_tF_j
    -\{J_4(u_j),F_j\}.
\end{multline*}
两式相减，得到
\begin{align*}
  \Delta H&-\operatorname{Re}(B_2-B_1)\partial_s^2F_1
    -\operatorname{Re}(B_2)\partial_s^2H
    +\operatorname{Im}(A_2-A_1)\partial_t^2F_1
    +\operatorname{Im}(A_2)\partial_t^2H\\
  &+\operatorname{Re}\bigl(A_2-A_1+i(B_2-B_1)\bigr)
       \partial_s\partial_tF_1
    +\operatorname{Re}(A_2+iB_2)\partial_s\partial_tH\\
  &=\operatorname{Re}(C_2-C_1)\partial_sF_1
     +\operatorname{Re}(C_2)\partial_sH
     +\operatorname{Im}(C_2-C_1)\partial_tF_1
     +\operatorname{Im}(C_2)\partial_tH\\
  &\quad-\{J_4(u_2)-J_4(u_1),F_1\}-\{J_4(u_2),H\}.
\end{align*}
令 $a:=1-\operatorname{Re}(B_2)$、$2b:=\operatorname{Re}(A_2+iB_2)$、$c:=1+\operatorname{Im}(A_2)$。于是
\[
  a\partial_s^2H+2b\partial_s\partial_tH+c\partial_t^2H
   =\lambda_1\cdot H+\lambda_2\cdot\partial_sH
                        +\lambda_3\cdot\partial_tH,
\]
其中 $\lambda_1\in\mathcal O_k$，$\lambda_2,\lambda_3\in\mathcal O_{k-1}$，取值于 $\operatorname{End}_{\mathbb R}(\mathbb C^{n-1})$。这些映射尤其有界，因此引理得证。
\end{proof}

\begin{proof}[命题~\ref{prop:E-6-1} 的证明]
在命题的记号下，对 $j=1,2$ 有 $U_j\subset\mathbb C$ 为零点邻域、$u_j:\overline U_j\to\mathbb C^n$ 为 \BookRecordIndex{T0003}$J$-全纯曲线，$\varphi_j:\mathbb D\to\overline U_j$ 为 $C^2$ 微分同胚，且 $\varphi_j-\mathrm{id}\in\mathcal O_2$；当 $k=1$ 时 $\varphi_j$ 光滑，并有
\[
  u_j\circ\varphi_j(z)=(z^k,F_j(z)),\qquad j=1,2,
\]
其中 $F_j\in\mathcal O_{k+1}$，$k=1$ 时 $F_j$ 光滑。设 $F_2-F_1$ 在零点的任意邻域内均不恒为零。须证明 $F_2-F_1=z^ma+R$，其中 $a\ne0$、$R\in\mathcal O_{m+1}$。

先设 $J$ 满足正规形条件~\eqref{eq:E-6-1}。对命题所给的 $F_1,F_2$，引理~\ref{lem:E-6-4}(iii) 表明 $C^2$ 函数 $H:=F_2-F_1$ 在零点邻域内满足 \BookRecordIndex{T0024}Hartman--Wintner 定理的假设，其中适当的 $a,b,c:\mathbb D\to\mathbb R$ 具有 Lipschitz 连续的一阶导数，并满足式~\eqref{eq:E-4-1}。由于 $H$ 在零点的任何邻域内都不恒为零，由定理~\ref{thm:E-4-1}，存在正整数 $m$ 及非零的 $m$ 次齐次多项式 $h:\mathbb C\to\mathbb C^{n-1}$，使
\[
  H(z)=h(z)+o(z^m).
\]
因 $H\in\mathcal O_{k+1}$ 且 $h\ne0$，有 $m>k$。还须证明 $h$ 全纯且 $H-h\in\mathcal O_{m+1}$。

设 $H\in\mathcal O_\ell$，其中 $\ell\le m$。当 $k>1$ 时，式~\eqref{eq:E-6-8} 给出 $\bar\partial H\in\mathcal O_{\ell+1}$；当 $k=1$ 时则给出 $\bar\partial H\in\mathcal O_\ell$。两种情形下，引理~\ref{lem:E-5-9} 都给出次数至多 $\ell$ 的复多项式 $p$，使
\[
  H-p\in\mathcal O_{\ell+1}.
\]
因 $\ell\le m$，$h-p=(h-H)+(H-p)=o(z^\ell)$。若 $\ell<m$，则 $p=o(z^\ell)$，故 $p=0$ 且 $H\in\mathcal O_{\ell+1}$。归纳得 $H\in\mathcal O_m$。取 $\ell=m$ 重复论证，由 $h-p=o(z^m)$ 得
\[
  h(z)=p(z)=az^m
\]
其中 $0\ne a\in\mathbb C^{n-1}$，并且 $H-az^m\in\mathcal O_{m+1}$。这就证明了正规形情形。

一般情形下，命题~\ref{prop:E-6-2} 给出局部光滑微分同胚 $\Phi:(\mathbb C^n,0)\to(\mathbb C^n,0)$，使
\[
  d\Phi(0)(\mathbb C\times\{0\})=\mathbb C\times\{0\},
  \qquad
  \widetilde J:=\Phi^*J
      =\begin{pmatrix}i&\widetilde J_3\\0&\widetilde J_4\end{pmatrix},
\]
其中 $\widetilde J_3(0)=0$、$\widetilde J_4(0)=i$；所以 $d\Phi(0)$ 具有式~\eqref{eq:E-6-10} 的形式。对 $j=1,2$，映射 $\widetilde u_j:=\Phi^{-1}\circ u_j:U_j\to\mathbb C^n$ 是 $\widetilde J$-全纯曲线。按假设及注~\ref{rem:E-5-11}，它在 $z=0$ 处的复切空间为 $\mathbb C\times\{0\}$，阶数为 $k$。由命题~\ref{prop:E-5-10}，存在局部 $C^{k,1}$ 微分同胚 $\widetilde\varphi_j:(\mathbb D,0)\to(U_j,0)$，使
\[
  \widetilde\varphi_j-\mathrm{id}\in\mathcal O_2,
  \qquad
  \widetilde u_j\circ\widetilde\varphi_j(z)
   =(z^k,\widetilde F_j(z)),
\]
其中 $k=1$ 时 $\widetilde\varphi_j$ 光滑，且 $\widetilde F_j\in\mathcal O_{k+1}$。由正规形情形，存在整数 $m>k$ 与非零向量 $\widetilde a\in\mathbb C^{n-1}$，使
\[
  \widetilde H-\widetilde p\in\mathcal O_{m+1},\qquad
  \widetilde H:=\widetilde F_2-\widetilde F_1,\quad
  \widetilde p(z):=\widetilde a z^m.
\]
在下述引理~\ref{lem:E-6-5} 中取 $v_j:=\widetilde u_j\circ\widetilde\varphi_j$、$G_j:=\widetilde F_j$、$\psi_j:=\widetilde\varphi_j^{-1}\circ\varphi_j$，即可得到
\[
  H-p\in\mathcal O_{m+1},\qquad
  H:=F_2-F_1,\quad p(z):=C\widetilde a z^m.
\]
命题得证。
\end{proof}

给定整数 $k$ 与映射 $f:\mathbb D\to\mathbb C^m$、$g:\mathbb D\to\mathbb C^\ell$，若存在 $\alpha:\mathbb D\setminus\{0\}\to\operatorname{Hom}_{\mathbb R}(\mathbb C^\ell,\mathbb C^m)$，$\alpha\in\mathcal O_k$，并在充分小的 $z$ 处满足 $f(z)=\alpha(z)\cdot g(z)$，就记 $f=\mathcal O_k\cdot g$。

\begin{lemma}\label{lem:E-6-5}
设整数 $k\ge1$，局部微分同胚 $\Phi:(\mathbb C^n,0)\to(\mathbb C^n,0)$ 满足
\begin{equation}\label{eq:E-6-10}
  d\Phi(0):=\begin{pmatrix}A&B\\0&C\end{pmatrix},
\end{equation}
其中 $A\in\mathbb C\setminus\{0\}$、$B\in\operatorname{Hom}_{\mathbb C}(\mathbb C^{n-1},\mathbb C)$、$C\in\operatorname{GL}(n-1,\mathbb C)$。对 $j=1,2$，设 $v_j:(\mathbb C,0)\to(\mathbb C^n,0)$、$\psi_j:(\mathbb C,0)\to(\mathbb C,0)$ 满足 $\psi_j-\mathrm{id}\in\mathcal O_2$，且
\[
  v_j(z)=(z^k,G_j(z)),\qquad
  \Phi\circ v_j\circ\psi_j(z)=(z^k,F_j(z)),
\]
其中 $F_j,G_j\in\mathcal O_{k+1}$。则对任意 $b\in\mathbb C^{n-1}$、$m\in\mathbb Z$，
\[
  G_2(z)-G_1(z)-bz^m\in\mathcal O_{m+1}
  \quad\Longleftrightarrow\quad
  F_2(z)-F_1(z)-Cbz^m\in\mathcal O_{m+1}.
\]
\end{lemma}

\begin{proof}
记 $\Phi_2$ 为 $\Phi$ 的后 $n-1$ 个坐标，$\Psi_1$ 为 $\Psi:=\Phi^{-1}$ 的首坐标。于是 $\Phi_2(\psi_j(z)^k,G_j(\psi_j(z)))=F_j(z)$，从而
\[
  H:=F_2-F_1
    =\Phi_2(\psi_2^k,G_2\circ\psi_2)
     -\Phi_2(\psi_1^k,G_1\circ\psi_1)=H_1+H_2,
\]
其中
\begin{align*}
  H_1&:=\Phi_2(\psi_2^k,G_1\circ\psi_2)
        -\Phi_2(\psi_1^k,G_1\circ\psi_1),\\
  H_2&:=\Phi_2(\psi_2^k,G_2\circ\psi_2)
        -\Phi_2(\psi_2^k,G_1\circ\psi_2).
\end{align*}
用 $d_1$ 表示关于 $\mathbb C^n=\mathbb C\times\mathbb C^{n-1}$ 首分量坐标 $z_1=s_1+it_1$ 的微分，用 $d_2$ 表示关于其余分量的微分。由注~\ref{rem:E-5-5}，映射 $z\mapsto d_1\Phi_2(z^k,G_j(z))$ 及 $z\mapsto d_2\Phi_2(z^k,G_j(z))-C$ 都属于 $\mathcal O_k$。因此，记 $K:=G_2-G_1$，有
\[
  H_2=\left(\int_0^1d_2\Phi_2
     (\psi_2^k,G_1\circ\psi_2+\tau K\circ\psi_2)\,d\tau\right)
       \cdot(K\circ\psi_2)
     =(C+\mathcal O_k)\cdot(K\circ\psi_2).
\]
再令 $\Gamma(z):=\Phi_2(z^k,G_1(z))$，则 $H_1=\Gamma\circ\psi_2-\Gamma\circ\psi_1$，并有
\[
  d\Gamma(z)=kz^{k-1}d_1\Phi_2(z^k,G_1(z))
              +d_2\Phi_2(z^k,G_1(z))\,dG_1(z)
  \in\mathcal O_k.
\]
因此
\[
  H_1=\left(\int_0^1d\Gamma\bigl(\psi_1+\tau(\psi_2-\psi_1)\bigr)
      \,d\tau\right)\cdot(\psi_2-\psi_1)
      =\mathcal O_k\cdot(\psi_2-\psi_1).
\]
另一方面，$\psi_j(z)^k=\Psi_1(z^k,F_j(z))$，故
\[
  \psi_2^k-\psi_1^k
    =\left(\int_0^1d_2\Psi_1(z^k,F_1+\tau H)\,d\tau\right)
      \cdot H=\mathcal O_0\cdot H.
\]
又因 $\psi_j-\mathrm{id}\in\mathcal O_2$，有
\[
  \psi_2^k-\psi_1^k
   =(kz^{k-1}+\mathcal O_k)(\psi_2-\psi_1)
   =kz^{k-1}(1+\mathcal O_1)(\psi_2-\psi_1).
\]
由上两式可得 $\psi_2-\psi_1=\mathcal O_{1-k}\cdot(\psi_2^k-\psi_1^k)=\mathcal O_{1-k}\cdot H$，从而 $H_1=\mathcal O_1\cdot H$。于是
\[
  H=H_1+H_2
   =\mathcal O_1\cdot H+(C+\mathcal O_k)\cdot K,
\]
故
\[
  H=(\mathbb 1+\mathcal O_1)^{-1}(C+\mathcal O_k)\cdot K
    =(C+\mathcal O_1)\cdot K.
\]
因此对任意 $b\in\mathbb C^{n-1}$ 和整数 $m$，
\[
  K(z)=z^mb+\mathcal O_{m+1}
  \quad\Longrightarrow\quad
  H(z)=z^mCb+\mathcal O_{m+1}.
\]
反向蕴涵只需以 $\Phi^{-1}$ 代替 $\Phi$，并以 $\Phi\circ v_j\circ\psi_j$ 代替 $v_j$。引理得证。
\end{proof}
